Das Seil verlässt die Trommel an ihrem Rand, also ist der Hebelarm der Last der Trommelradius $r$; der Hebelarm der Hand ist die Kurbellänge $R$. Im Gleichgewicht sind die beiden Drehmomente bezüglich der Achse gleich gross:
\begin{align}
 F\cdot R &= M\,g\cdot r
\end{align}
Also
\begin{align}
 F &= \FKF \\
   &= \frac{\ML\cdot\g\cdot\rT}{\RK} \\
   &= \FK \approx \result{\FKP{0}{2}}
\end{align}
Ein Fünftel der Gewichtskraft der Last: Die Kurbel ist fünfmal so lang wie der Trommelradius.
